\scanpage{188}\mspage{44}
% archive sheet number: 188

est iso.

$\deleted{\unsure{b}}a')$ Pour tte \deleted{\unsure{cat}} $\deleted{\unsure{M}}\mathcal{U}$-catégorie $M$ où les petites limites projectives sont représentables, \deleted{la flèche} \inserted{et tout $F\in M(I)=\underline{\mathrm{Hom}}(I^\circ,M)$}, \deleted{\unsure{pr} tout} \unsure{préfaisceau} sur $I$ à val. dans $M$, la flèche can\unsure{.}

\[
\underset{I}{\underset{\longleftarrow}{\lim}}\,F\longrightarrow\underset{I'}{\underset{\longleftarrow}{\lim}}\,u^*(\deleted{\unsure{G}}F)\tag*{(3.2.2.)}
\]

est iso.

$\deleted{\unsure{b}}a'')$ Comme a'), avec $M=(\mathrm{Ens})$ i.e. $F\in\mathrm{Ob}\,I^\wedge$.

% slanted
\margin{C'est une condition sur des morphismes de topos, savoir que $\forall F$ \deleted{dans} $\unsure{\mathscr{X}|X}$, $H^0(\mathscr{X},F)\overset{\sim}{\longrightarrow}H^0(\unsure{\mathscr{X}'},u^*F)$. Il y a notion similaire pour les $H^i$ $(i\leq n)$, \unsure{à dégager|à dég.} Notion des morphismes $W_0$-coasphériques des \unsure{catégories|objets} \unsure{fibrés|fléchés} d'un topos ! C'est \unsure{donc|dans} la \unsure{notion de dualité|notion de validité}, condition de \unsure{$g$|$G$}, \unsure{ou|vu} \unsure{comme} $W_0$-\unsure{asphérique}, qui \unsure{dérive|doute} de la notion \unsure{de} $\unsure{U^\wedge|U'}$ \unsure{induit} \unsure{des} signifie simplement que pour \unsure{tout} \unsure{objet} \unsure{dont le} $U^\wedge$, $\unsure{f^*|T^*}(U)$ est \unsure{$0$-connexe|connexe}. \unsure{Ici} \unsure{il convient de} \unsure{comparer} une notion \unsure{anciennement|``ancien''} de \deleted{\illegible{1}} $W_\infty$-asphéricité et \unsure{bien comprise}, \uline{\unsure{précisément}} \unsure{(ancien} $u$-asphérique), et (\deleted{\unsure{donc}} $W_\infty$-\unsure{coasphéricité}) \deleted{\unsure{est}} loc. \unsure{3.2.3.} \unsure{étant} celle évoquée plus haut \unsure{a''}.}

\margin{\uline{NB} On peut \unsure{intercaler|encore avoir} $a''')$ Comme a), mais avec $M=\mathrm{Ens}$ $a^{\mathrm{iv}})$ \deleted{\unsure{Comme}} a) avec le foncteur canonique $I'\longrightarrow\unsure{I'{}^\wedge|I^\wedge}$.}

b) La flèche $u$ est $W_0$-coasphérique, i.e. pour tout $i\in I$, $\,{}_{i\backslash}I'$ est $0$-connexe.

L'équivalence de a, a', a'' est $\pm$ tautologique, \unsure{compte tenu|vu} des définitions (cf \unsure{SGA 4 I 8.1}).

La a'') s'écrit, posant $\mathbb{D}=\mathrm{HOT}_{W_0}$, $\forall F\in\mathbb{D}(I)$, on veut que

\[
\begin{aligned}
H^{\unsure{\bullet|0|\cdot}}_{\mathbb D}(I,F)&\longrightarrow H^{\bullet}_{\mathbb D}(I',\deleted{\unsure{M}}u^*F)\\
\Vert\ (\mathrm{déf})&\qquad\Vert\ \mathrm{déf}\\
p_{I*}(F)&\qquad p_{I'*}(F)
\end{aligned}\tag*{(3.2.3.)}
\]

($p_I:I\longrightarrow e$ et $p_{I'}:I'\longrightarrow e$ les foncteurs canoniques) soit iso. On sait que cel\unsure{a} équivaut à $u$ $\mathbb{D}$-coasphérique. La démonstration est un jeu amusant entre $u_!$ \unsure{et} $u^*$, $u_*$ à la fin. Il faut exprimer que le $\mathrm{Hom}$ d'un objet $M$ de $\mathbb{D}(e)$ ($=\mathrm{Ens}$ en l'occurrence) vers la flèche (3.2.3) est bijective. Ici, il
